\section{本章小结}
我们从手写规则走到数据学习，先用独立验证点区分拟合与泛化，再把图片与标签分成训练、验证和测试。卷积、池化与激活组成特征，展平后用全连接输出 logits；损失给出差距，反向传播求梯度，优化器才更新参数。训练先用少量数据检查闭环，再恢复完整集合，按验证损失保存模型，最后才评价测试集。

你应能熟练核对输入、每层形状、设备与训练状态，保存结构、预处理和类别说明，重新载入指定样本并解释分数。检测还需要框解码与后处理，换运行时要核对完整数值，再测用时；现成检测框架、OpenVINO/TensorRT 的具体支持和赛场数据继续查对应版本与卷三，不能凭文件后缀或算力推断效果。

\begin{problemset}
\item\label{prob:v2-dl-1} 基础：训练误差下降而验证误差上升，优先检查什么？同一段视频的相邻帧被随机分到训练与验证，会怎样影响评价？
\item\label{prob:v2-dl-2} 基础：输入 $3\times10\times10$、输出 8 通道、$3\times3$ 核、每侧填充 1、步幅 2 的卷积，输出形状与参数数目（含偏置）是多少？
\item\label{prob:v2-dl-3} 基础：28 输入的本章网络为什么有 256 个展平特征？把 batch 从 1 改成 8，需要将 Linear 输入数乘以 8 吗？
\item\label{prob:v2-dl-4} 基础：logits、softmax 概率与交叉熵是什么关系？backward 后参数是否已更新，第二次 backward 是否自动清梯度？
\item\label{prob:v2-dl-5} 基础：eval 与 \texttt{no\_grad} 分别控制什么？本章 LeNet 没有 Dropout/BatchNorm，是否就可以永远不写 eval？
\item\label{prob:v2-dl-6} 基础：交付 best.pth 还需要哪些说明？输入形状一致是否足以证明两个预处理流程等效？
\item\label{prob:v2-dl-7} 提高：驱动可用但 \texttt{torch.cuda.is\_available} 为 False，应先检查什么？GPU 前向只测普通时钟间隔为什么可能低估时间？
\item\label{prob:v2-dl-8} 提高：历史 YOLOv5 的 255 个预测通道代表什么？改成 3 类时该通道数是多少？同类 NMS 能否不加区分地用于不同类？
\item\label{prob:v2-dl-9} 提高：本章 ONNX 文件能直接输入批量 8 吗？导出检查通过、某张图类别相同，是否足以宣布部署一致？
\item\label{prob:v2-dl-10} 实践：独立加载训练选中的参数，对检查者指定的三个测试索引输出分数；其中一个再跨运行时比较。提交结构、预处理、类别、误差与 CPU 前向计时，并说明实机单帧响应还缺什么。
\end{problemset}

\section{本章验收}
\begin{enumerate}
  \item 从独立环境运行小规模与完整 MNIST 训练，展示版本、首批形状、三个集合大小、逐轮曲线及最终测试；说明参数选择没有使用测试集。
  \item 现场改变卷积/池化步幅或输入尺寸，手算并验证输出，解释维度错误；展示梯度计算、累加、清零与更新的区别。
  \item 在另一进程加载最佳参数，对指定索引推理，解释预处理与类别表；有 CUDA 时展示同参数输出比较，无设备如实记录缺口。
  \item 对同一真实输入比较 PyTorch 与 ONNX Runtime 完整分数，检查输入限制并展示 CPU 批量计时，区分推理成本、吞吐与相机端到端响应。
\end{enumerate}
